% from main.tex
%%%%%%%%%%  scan idx 421  |  folio 408  %%%%%%%%%%
homomorphisme d'augmentation : $K^{\bullet}(X)\longrightarrow H^0(X,\mathbb Z)$.

\textbf{Proposition 6.1.} Soit $X$ un schéma quasi-compact ayant un faisceau ample $L$. Alors $\Fil^1K^{\bullet}(X)$ est un nilidéal.

Comme $\Fil^1K^{\bullet}(X)$ est engendré par les éléments du type $E-\varepsilon(E)$, où $E$ est la classe d'un $\mathcal O_X$-Module localement libre de type fini $E$, et $\varepsilon:K^{\bullet}(X)\to H^0(X,\mathbb Z)$ l'augmentation de $K^{\bullet}(X)$, il suffit de montrer qu'un tel élément est nilpotent. La quasi-compacité de $X$, et la commutation du $K^{\bullet}$ aux sommes directes finies permet de se ramener au cas où $E$ est de rang constant. Enfin, en passant au fibré de drapeaux de $E$, ce qui donne une injection sur les $K^{\bullet}$ (3.1), et conserve l'hypothèse d'éxistence d'un faisceau ample, on peut supposer $E$ inversible.

Soit $E^*$ le dual de $E$. Pour $n$ assez grand, $E^*(n)=E^*\otimes_{\mathcal O_X}L^{\otimes n}$ est engendré par ses sections globales; on peut donc trouver un entier $m$, et un homomorphisme surjectif : $(\mathcal O_X)^m\longrightarrow E^*(n)$, donc (tensorisant par $E(-n)$) un homomorphisme surjectif : $(E(-n))^m\longrightarrow\mathcal O_X$. Le noyau de cet homomorphisme est localement libre de rang $m-1$. On obtient donc, en notant $L$ la classe de $L$:
\[
  \lambda^m(mEL^{-n}-1)=0 .
\]
Or $\lambda^m(mEL^{-n}-1)=(-1)^m\lambda_{-1}(mEL^{-n})=(-1)^m\lambda_{-1}(EL^{-n})^m=(-1)^m(1-EL^{-n})^m$. Par conséquent, pour tout faisceau inversible $E$, l'élément $1-EL^{-n}$ est nilpotent pour $n$ assez grand, ou encore son produit $L^n-E$ par $L^n$ est nilpotent.

On peut alors écrire:
\[
  E-1=(L^n-1)-(L^n-E)
\]
et d'après ce qui précède, $L^n-E$ et $L^n-1$ sont nilpotents si on choisit $n$ assez grand. Donc $E-1$ est nilpotent.


